\chapter{Quick Reference}\label{ch:reference}

\chapterguide%
  {Everything in this appendix has appeared earlier with full explanation; nothing here is new material.}
  {revise before an assessment, choose a structure for the final project, and check a complexity claim quickly.}

This appendix collects the numbers and decisions scattered across the previous chapters into
one place. Use it for revision, not for first learning: each entry points back to the chapter
that explains it.

\section{Choosing a data structure}

Most structure choices in this course reduce to a short series of questions about the access
pattern (\figref{fig:choose}). The pattern matters more than the data: the same student
records belong in an array, a BST, or a hash table depending entirely on how they will be
looked up.

\begin{adsfig}[Choosing a data structure]{A decision path for the structures covered in this course. Read it as a set of
questions about how the data will be \emph{used}, not about what the data
is.}{fig:choose}
\begin{tikzpicture}[font=\small,
  q/.style={draw=ADSOrange,fill=ADSPaleOrange,rounded corners=2mm,text width=3.5cm,
            align=center,inner sep=2.2mm,font=\scriptsize\color{ADSNavy}},
  ans/.style={draw=ADSGreen,fill=ADSPaleTeal,rounded corners=2mm,text width=2.8cm,
            align=center,inner sep=2.2mm,font=\scriptsize\bfseries\color{ADSNavy}},
  yes/.style={adsarrow,draw=ADSGreen},
  no/.style={adsarrow,draw=ADSRed}]

\node[q] (q1) at (0,0) {Do you look items up by a \textbf{key}?};
\node[q] (q2) at (0,-2.4) {Is order (sorted, LIFO, FIFO) part of the requirement?};
\node[q] (q3) at (0,-4.8) {Do items have \textbf{hierarchical or network} relationships?};
\node[ans] (a0) at (5.5,0) {Hash table\\{\scriptsize\mdseries\chapref{ch:hashing}}};
\node[ans] (a1) at (5.5,-2.4) {Stack / Queue / BST\\{\scriptsize\mdseries Ch.~\ref{ch:stacks}, \ref{ch:queues}, \ref{ch:trees-ii}}};
\node[ans] (a2) at (5.5,-4.8) {Tree / Graph\\{\scriptsize\mdseries Ch.~\ref{ch:trees-i}, \ref{ch:graphs}}};
\node[ans] (a3) at (0,-7.0) {Array or linked list\\{\scriptsize\mdseries Ch.~\ref{ch:arrays}, \ref{ch:linked-lists}}};

\draw[yes] (q1)--(a0) node[midway,above,adslabel,color=ADSGreen] {yes};
\draw[yes] (q2)--(a1) node[midway,above,adslabel,color=ADSGreen] {yes};
\draw[yes] (q3)--(a2) node[midway,above,adslabel,color=ADSGreen] {yes};
\draw[no] (q1)--(q2) node[midway,left,adslabel,color=ADSRed] {no};
\draw[no] (q2)--(q3) node[midway,left,adslabel,color=ADSRed] {no};
\draw[no] (q3)--(a3) node[midway,left,adslabel,color=ADSRed] {no};

\node[adslabel,anchor=west,align=left,text width=5.2cm] at (7.3,-7.0)
  {Then choose between them on one further question: is the size fixed and is
   direct index access needed (array), or will there be many insertions and deletions
   (linked list)?};
\end{tikzpicture}
\end{adsfig}

\section{Operation costs by structure}

\begin{mltable}
\begin{tabularx}{\linewidth}{@{}>{\bfseries\leavevmode\color{ADSNavy}}p{0.20\linewidth}
  >{\centering\arraybackslash}p{0.135\linewidth}
  >{\centering\arraybackslash}p{0.135\linewidth}
  >{\centering\arraybackslash}p{0.135\linewidth}X@{}}
\toprule
\rowcolor{ADSPaleBlue!92}
Structure & Access & Search & Insert/Delete & Note \\
\midrule
Array                  & $O(1)$      & $O(n)$      & $O(n)$      & Shifting dominates insertion \\
Sorted array           & $O(1)$      & $O(\log n)$ & $O(n)$      & Binary search applies \\
Singly linked list     & $O(n)$      & $O(n)$      & $O(1)$      & Cost is $O(1)$ \emph{once positioned} \\
Doubly linked list     & $O(n)$      & $O(n)$      & $O(1)$      & Deletion needs no predecessor search \\
Stack                  & $O(1)$ top  & ---         & $O(1)$      & Access restricted to the top \\
Queue                  & $O(1)$ ends & ---         & $O(1)$      & Circular form reuses slots \\
BST (balanced)         & ---         & $O(\log n)$ & $O(\log n)$ & Degrades to $O(n)$ if unbalanced \\
Heap                   & $O(1)$ root & $O(n)$      & $O(\log n)$ & Root is min or max only \\
Hash table             & ---         & $O(1)$ avg  & $O(1)$ avg  & Worst case $O(n)$ under heavy collision \\
\bottomrule
\end{tabularx}
\end{mltable}

\begin{clarification}
These are standard averages for the implementations built in this course, not figures quoted
from the slide decks. Where a chapter states a different value, the chapter's own
clarification box explains the discrepancy.
\end{clarification}

\section{Sorting algorithms at a glance}

\begin{mltable}
\begin{tabularx}{\linewidth}{@{}>{\bfseries\leavevmode\color{ADSNavy}}p{0.17\linewidth}
  >{\centering\arraybackslash}p{0.15\linewidth}
  >{\centering\arraybackslash}p{0.15\linewidth}
  >{\centering\arraybackslash}p{0.10\linewidth}X@{}}
\toprule
\rowcolor{ADSPaleBlue!92}
Algorithm & Average & Worst & Stable & Choose it when \\
\midrule
Bubble    & $O(n^2)$     & $O(n^2)$     & yes & Teaching only; rarely practical \\
Insertion & $O(n^2)$     & $O(n^2)$     & yes & Small or nearly-sorted input \\
Selection & $O(n^2)$     & $O(n^2)$     & no  & Swaps are expensive; only $n-1$ occur \\
Merge     & $O(n\log n)$ & $O(n\log n)$ & yes & Guaranteed bound and stability needed \\
Quick     & $O(n\log n)$ & $O(n^2)$     & no  & Fast in practice; in-place \\
Heap      & $O(n\log n)$ & $O(n\log n)$ & no  & Guaranteed bound with $O(1)$ extra space \\
\bottomrule
\end{tabularx}
\end{mltable}

\clearpage
\section{Traversal orders}

\begin{mltable}
\begin{tabularx}{\linewidth}{@{}>{\bfseries\leavevmode\color{ADSNavy}}p{0.20\linewidth}p{0.30\linewidth}X@{}}
\toprule
\rowcolor{ADSPaleBlue!92}
Traversal & Rule & Typical use \\
\midrule
Inorder   & left, root, right & Sorted output from a BST \\
Preorder  & root, left, right & Copying or serialising a tree \\
Postorder & left, right, root & Deleting or freeing a tree safely \\
BFS       & queue, level by level & Shortest path in an unweighted graph \\
DFS       & stack or recursion & Cycle detection, connectivity, backtracking \\
\bottomrule
\end{tabularx}
\end{mltable}

\section{Pointer operations that go wrong most often}

\begin{mltable}
\begin{tabularx}{\linewidth}{@{}>{\bfseries\leavevmode\color{ADSNavy}}p{0.28\linewidth}X@{}}
\toprule
\rowcolor{ADSPaleBlue!92}
Mistake & What to do instead \\
\midrule
Moving \texttt{head} before linking the new node & Set \texttt{n->next = head} first (\figref{fig:sll-insert-head}) \\
Deleting a node before relinking its predecessor & Redirect \texttt{prev->next}, then \texttt{delete} (\figref{fig:sll-delete}) \\
Using \texttt{while (p != nullptr)} on a circular list & Use \texttt{do}\,--\,\texttt{while} against the head (\figref{fig:cll-traverse}) \\
Popping or peeking without an empty check & Test \texttt{isEmpty()} before touching the top \\
Treating a full linear queue as genuinely full & Use modular arithmetic (\figref{fig:queue-circular}) \\
Omitting the visited set in graph traversal & Mark on discovery, not on processing (\figref{fig:visited}) \\
\bottomrule
\end{tabularx}
\end{mltable}
